\subsection{Transformers and Large-Scale Training}

Transformers combine attention, feedforward blocks, normalization, and residual connections into architectures that scale well on sequence and multimodal data.

\begin{definition} \label{def:sl-transformer-block} \hlt{Transformer Block}\\
A transformer block applies multi-head self-attention, residual connections, normalization layers, and a positionwise feedforward network.  Stacking these blocks yields deep contextual representations.
\end{definition}

\begin{definition} \label{def:sl-self-attention} \hlt{Self-Attention}\\
Self-attention uses the same sequence to form queries, keys, and values.  Multi-head attention runs several attention maps in parallel and recombines their outputs.
\end{definition}

\begin{remark} \label{rem:sl-large-scale-training} \hlt{Large-Scale Training}\\
Large-scale deep learning is constrained by optimization, data quality, memory, communication, and evaluation.  Batch size, learning-rate schedules, mixed precision, parallelism, and checkpointing are part of the mathematical and systems design of training.
\end{remark}

\begin{remark} \label{rem:sl-transformer-generalization} \hlt{Inductive Bias}\\
Transformers have weaker built-in locality bias than convolutional networks but stronger capacity for long-range interaction.  Domain-specific variants add locality, recurrence, sparsity, or equivariance when the task demands it.
\end{remark}
